% Exercice à caractère expérimental
L'hydroxyde de sodium, plus communément appelé soude, se présente sous la forme
de pastilles blanches.

\begin{tcolorbox}[sidebyside,righthand width=2cm,lower separated=false]
    Le pictogramme figurant sur un flacon est représenté ci-contre. En outre, il
    est inscrit sur l'étiquette du flacon~: <<R:25 et S/26-37/39-45>>.  
    \tcblower

    \begin{figure}[H]
        \centering
        \includegraphics[width=\textwidth]{2025-12-09_011436-}
    \end{figure}
\end{tcolorbox}

\begin{enumerate}
    \item Quelles précisions sur les dangers de cette substance le pictogramme
          fournit-il~?
    
          Rechercher dans un catalogue de produits chimiques les risques $(R)$
          et consignes de sécurité $(S)$ attachés à l'hydroxyde de sodium.
    \item On veut préparer $\qty{250}{\mL}$ d'une solution d'hydroxyde de sodium
          de concentration $\qty{0.30}{\mol\per\L}$. Décrire
          le mode opératoire permettant d'effectuer cette préparation.
    \item Écrire l'équation de la réaction de dissolution de l'hydroxyde de
          sodium dans l'eau~; en déduire la concentration molaire des ions
          présents dans la solution.
    \item Sachant que l'on ne peut fractionner une pastille et que la masse
          d'une pastille est de $\qty{0.09}{\g}$, quelle concentration la plus
          voisine de celle désirée pourra-t-on obtenir~?
\end{enumerate}

